\section{Introduction}

A coding agent is given standing permission to act: edit files, run builds,
commit. It still stops early. It names the next step instead of taking it,
offers a menu instead of picking, reports a result and asks whether to
continue, or hands the user a command the agent could have run itself. The
user's reply is then one word. That word costs a round trip, breaks the
user's attention, and carries no information the agent did not already have.

The failure is not a capability failure. The same agent, woken with nothing
more than ``keep going,'' usually finishes the work. It is a stopping-policy
failure, and it is invisible to the agent, which believes each turn ended
reasonably.

This paper describes a gate that sits in the agent's stop path and re-wakes it
when the closing message shows the pattern. That part is straightforward. The
part that took the work, and the part we think transfers, is the second loop:
how the gate discovers which of its own interventions were wrong, from signals
already present in the transcript, with no human labeling.

We contribute:

\begin{enumerate}
  \item A three-lane gate that spends a regular expression before it spends a
        model, and a model before it spends the user's attention
        (Section~\ref{sec:system}).
  \item A weak-supervision label for premature stopping that requires no
        annotation: the user's own next message (Section~\ref{sec:setup}).
  \item An evaluation over 747 clean closings from 43 sessions, including an
        ablation of every lane and every pattern class against that label (Section~\ref{sec:results}).
  \item A self-labeling loop that grades each intervention by what the agent
        did afterwards, and names patterns for demotion on its own
        (Section~\ref{sec:loop}).
  \item Negative results, in full, including two failed attempts to obtain a
        usable confidence signal from the judge (Section~\ref{sec:negative}).
\end{enumerate}

\section{Background and related work}

\paragraph{Where the gate runs.} Claude Code exposes a \emph{Stop hook}: a
program invoked when the agent finishes a turn. Exit status 0 lets the turn
end. Exit status 2 discards the ending and wakes the agent with the hook's
message on standard error. The hook sees the full transcript path, the working
directory, and a flag marking whether it is already inside a re-wake chain. It
is, in effect, an interposition point on the agent's stopping decision.

\paragraph{Model-as-judge.} Using one language model to score another's output
is now standard practice, and its failure modes are documented: position bias,
verbosity bias, and self-preference \cite{zheng2023judging}. Our judge is
smaller than the agent it judges by roughly two orders of magnitude, which
removes self-preference and forces the prompt to carry the whole task.

\paragraph{Self-critique and refinement.} Self-Refine \cite{madaan2023selfrefine}
and Reflexion \cite{shinn2023reflexion} have a model improve its own output
from its own feedback. Constitutional AI \cite{bai2022constitutional} supplies
the critique from written principles. Our setting differs in one respect that
matters: the agent has already decided it is done, so a self-critique would be
asking the party that made the error to detect it. The gate is deliberately
external.

\paragraph{Process supervision.} Rewarding intermediate steps rather than final
answers improves reliability \cite{lightman2023verify}. Premature stopping is a
process defect with a correct final answer -- the work the agent did is fine --
so it is invisible to outcome supervision.

\paragraph{Weak supervision.} Snorkel \cite{ratner2017snorkel} builds training
labels from noisy programmatic labeling functions rather than annotators. Our
pattern lane is a set of labeling functions in exactly that sense, and
Section~\ref{sec:loop} adds the piece Snorkel assumes you have: an estimate of
each function's accuracy, obtained here from the agent's later behaviour.

\paragraph{Calibration.} Language models carry some signal about their own
correctness \cite{kadavath2022know}, and self-consistency across samples is a
usable proxy \cite{wang2022selfconsistency}. Section~\ref{sec:negative} reports
that both failed here at 9B scale, and that token log-probabilities did not.

\section{System}
\label{sec:system}

The gate is eight checkers and one local model behind a single orchestrator:
5389 lines of Python and 18 test files. Three lanes decide, cheapest first, and
a lane that decides ends the turn's evaluation.

\subsection{Lane 0: auditing the escape hatch}

A closing message that begins with \texttt{BLOCKED:} is exempt from every other
check. An agent that cannot proceed must be able to say so, and a gate that
punishes a genuine blocker teaches the agent to lie about its state.

That exemption is also the cheapest escape in the system, and it was used as
one. Lane 0 therefore audits the claim against the transcript rather than the
prose. Three signals fail the audit: no tool call at all in the turn that
declared the blocker, so nothing was tried; four or more distinct blockers named
at once, which is a survey rather than a wall; and three consecutive blocked
closings, which is a loop. A blocker that survives the audit passes untouched.

An agent that survives the audit and still needs out is given a release phrase
drawn from a cryptographic random source, generated per block and present in no
file the agent can read. Quoting it back proves the agent read the current
block, and the blocker it accompanies is then audited against what actually ran
that turn. The phrase exists because every fixed escape token we tried was
eventually emitted from memory rather than from the message.

\subsection{Lane 1: patterns as labeling functions}

The second lane is 299 regular expressions over the agent's own words, grouped
into nine classes by what the closing is doing rather than by what it says:
\texttt{ask}, \texttt{announce}, \texttt{defer}, \texttt{wait},
\texttt{handoff}, \texttt{errand}, \texttt{read}, \texttt{idle} and
\texttt{ack}. Appendix~\ref{app:classes} gives the full table. The grouping
matters because the loop of Section~\ref{sec:loop} tunes classes rather than
individual expressions: one regular expression rarely fires often enough to be
judged on its own.

Two design points matter more than the patterns themselves.

\paragraph{Quoted text does not count.} Fenced code blocks, markdown tables,
inline backticks, quotation marks and reported speech are stripped before
matching. Without this, an agent that describes the gate trips it, and a report
containing a table of results is condemned by its own table. This was not
hypothetical: it fired on this project's own status reports twice before the
table rule was added.

\paragraph{Patterns carry a confidence.} A class written \texttt{ask:} decides
on its own. A class written \texttt{ask?:} raises a doubt and defers to the
model. Of 299 patterns, 222 are firm and 77 are doubts. This split is what makes
a local model affordable: most closings never reach it, and the ones that do are
the ones a regular expression was going to get wrong.

\subsection{Lane 2: a small local judge}

The third lane is a 9B parameter model served locally through Ollama, pinned by
version, answering with exactly one token: \textsc{stop} or \textsc{ok}. It runs
at temperature 0 with a fixed seed.

The judge exists for one distinction a regular expression cannot make. In this
corpus the Spanish word \emph{falta} appears both as ``work is missing'' and as
``it was missing and I fixed it.'' No amount of pattern engineering separates
those; reading the sentence does.

The judge reports how much probability mass it put on the token it chose, read
from the logprobs that accompany every response. That number was recorded from
the first verdict and used for nothing for two weeks. It separates cleanly: over
the 17 held-out messages the judge labels correctly, its lowest score is 0.846,
and every block in the production log scores above 0.5. The one verdict below
that threshold was a false accusation, scoring 0.228 on a message that closed
cleanly. A \textsc{stop} below 0.5 is therefore not delivered as a block; it is
delivered as the proactive question of Section~\ref{sec:proactive}, which buys
the same second look without telling the agent it left work behind. Guessing
costs a look; accusing costs trust.

A judge that cannot be reached is treated the same way. An earlier version
blocked every stop when the model was silent, on the argument that a switched-off
judge is precisely the failure the gate exists to prevent, and it started the
daemon itself to avoid that state. Both halves were wrong in the same direction.
Starting a 5.5GB model server because a turn ended is a decision belonging to
whoever owns the machine: on 2026-08-29 the gate did it on a busy workstation,
which ran out of memory and had to be restarted. The gate now starts nothing,
and a silent judge degrades to the pattern lanes plus one proactive question.

Two further implementation details proved necessary rather than cosmetic. First, all
judge requests are serialized behind a machine-wide lock: the serving daemon
batches concurrent requests, and batching changes the answer, so two test suites
running at once produced different verdicts on identical input. Second, when the
judge is asked to quote the offending sentence back to the agent, the prompt is
strictly extractive and forbids translation; without that, the model paraphrased
the sentence into English and the quote failed validation.

\subsection{The last question before the turn closes}
\label{sec:proactive}

An \textsc{ok} from the judge used to end the turn. It no longer does, when the transcript
can name what is missing. The gate reads the turn's own tool calls, collects
the files it wrote, and asks whether there is a next action available right
now, naming the candidates it found: the source file changed with no test
written beside it, the code changed with nothing in the documentation moving
with it. The turn ends when the agent answers that it looked and found
nothing, or when it goes and does the thing it found.

The candidates are read off the record rather than inferred, and a turn that
leaves none is released without a question. That gate on the question is
measured rather than assumed: a candidate exists on 30\% of real turns, and
the first blocks of the unconditional version bought no work half the time,
the worst conversion of any lane (Section~\ref{sec:results}).

The asymmetry is deliberate. A block says the agent left work behind, and it is
wrong whenever the closing was clean. This question makes no such claim, so it
costs one paragraph when the answer is no and buys a whole unit of work when
the answer is yes. It is asked once per closing: the state file records that the
question was put, and the answer to it is released without a second look.

That property is what makes it the right home for two other cases. A block the
judge scores below its confidence floor arrives here instead of as an
accusation, and so does a stop the judge never answered because it was switched
off. In both the gate is uncertain, and a question is what uncertainty is
entitled to.

\subsection{The message the agent receives}
\label{sec:tone}

A block writes one short message to the agent. Its wording went through several
revisions and the final form is deliberate: it states that the work already done
stands, that permission was granted in advance and does not expire, and that if
anything remains the agent should take the step rather than name it.

The reasoning is mechanical rather than courteous. An agent that reads an
accusation spends the woken turn defending itself, which is precisely the
outcome the block was meant to prevent. We observed this directly with earlier,
blunter wording.

The same reasoning keeps the diagnostics out of the message. Which patterns
fired, how confident the judge was, whether the lanes disagreed -- all of it is
useful for tuning and all of it reads as prosecution evidence. It is written to
a private log instead. The agent receives the request; the maintainer receives
the numbers. Section~\ref{sec:loop} is built on that log.

\section{Experimental setup}
\label{sec:setup}

\subsection{Corpus}

We use every session transcript on one developer's machine: 66 transcripts
spanning 2026-08-03 to 2026-08-29, across several projects. From these we
extract \emph{closing pairs}: the agent's final text message of a turn together
with the user's next message. The first build of this yielded 854 pairs. The
corrected readers of Section~\ref{sec:contamination} yield 764, of which 747
are neither harness noise nor the gate talking to itself.

The corpus is naturally split by a date. The gate was first installed on
2026-08-28, so of the 43 sessions that carry a closing, 32 are \emph{ungated}
- produced with no gate in the stop path - and 11 are \emph{gated}.
Ungated transcripts are the ones the gate never influenced. The lane
ablations run over the whole clean corpus, because a gated closing is still a
closing the developer answered; what the split is for is the era comparison,
and for reading the gate's misses, which is only honest where the gate was
not shaping the text.

\subsection{A label nobody had to write}

Evaluating a stop gate requires knowing which stops were premature. Annotating
that by hand is expensive and, worse, is annotated by the same person whose
intuitions built the patterns.

The corpus already contains the judgement. When an agent stops early, the user's
next message is a bare push: \emph{dale}, \emph{seguí}, \emph{continue},
\emph{keep going}. It carries no new instruction, because none was needed. We
detect it with a closed lexicon anchored to the whole message, so that
``continue with the refactor'' is not counted and ``ok dale'' is.

This label is noisy in one direction and clean in the other. A push is strong
evidence the turn ended too early. Its absence is weak evidence of the opposite,
because a user often absorbs the loss and just answers the question. We
therefore report recall against pushes and firing rate, and do not report
precision, which this label cannot support.

\subsection{Metrics}

For each lane configuration we report how many closings it fires on, and what
fraction of pushed closings it catches. We deliberately avoid a single combined
score: the two costs are not commensurable. A missed premature stop costs the
user one round trip. A false block costs the agent a full turn and, if the
wording is wrong, poisons the next one.

\section{The self-labeling loop}
\label{sec:loop}

Every decision the gate makes is appended to a private log: the lane that
decided, the patterns that fired, the judge's verdict and confidence, the
latency, and the first line of the message. The log records what was decided. It
does not record whether the decision was right, and for the first weeks tuning
meant reading transcripts by hand.

The missing label was already in the transcript, a few lines below the block.
A block wakes the agent, and from there exactly one of two things happens: the
agent goes and does work, or it writes another paragraph and stops again.

\begin{quote}
\textbf{Grading rule.} The number of tool calls between a block and the next
time the user speaks is the label for that block. Above a threshold, the block
bought work. At or below it, the block bought a paragraph.
\end{quote}

Nobody supplies this. It is read from the same transcripts the gate already
opens. The report joins each logged decision back to its transcript, counts the
intervening tool calls, and aggregates by pattern class and by judge confidence.
When a class accumulates enough blocks and most of them bought nothing, the
report names that class for demotion from firm to doubt.

The report runs at session start and stays silent unless a class is ripe,
speaking at most once per day. A report that speaks on every session is a report
that gets skipped, and the loop's output is consumed by an agent with the same
attention economics as its user.

The first two blocks the loop graded had bought zero tool calls each, on the
same pattern class that had been demoted by hand that morning after measuring
888 closings. The loop reproduced a hand measurement in one second.

\section{Negative results}
\label{sec:negative}

These cost more than the positive ones and generalize better.

\paragraph{The judge cannot report its own confidence.} Asked to append a
confidence rating to its verdict, the 9B model answered with the maximum every
time, across every message we tried. The number was constant and therefore
carried nothing.

\paragraph{Self-consistency did not separate either.} Re-sampling the same
message five times at temperature 0.8 produced five identical verdicts,
including on messages the verdict got wrong. At this scale the model is
confidently wrong in a stable way, so disagreement-based confidence
\cite{wang2022selfconsistency} has nothing to measure.

\paragraph{Token log-probabilities did move.} Taking the probability the model
assigned to the single verdict token it emitted produced a number that varies
with the message, and that separates at the extremes. It is the only one of the
three that worked.

\paragraph{A 9B judge cannot hold an exception.} Nineteen percent of the judge's
blocks were turns in which the agent had correctly launched background work and
was waiting on it. We tried to teach this in the prompt, first as prose and then
as few-shot examples. Both failed: the model could not hold the exception
against its own standing instruction to answer \textsc{stop} when unsure, and
the examples degraded unrelated verdicts. The fix was to remove the question.
Waiting is now established deterministically from the transcript, before the
judge is called at all. \emph{An instruction a small model cannot hold is a
decision that belongs outside the model.}

\paragraph{Nor a disposition.} The judge's instruction ended with
\emph{when unsure, answer STOP}. Replacing it with its exact opposite,
\emph{when unsure, answer OK}, changed 31 verdicts out of 841 and moved
precision by 0.001. Replacing it instead with a test the model can apply --
is there an action available right now that needs nothing only the developer
knows -- moved precision, recall and intervention count together. Three
instructions in this work went unheld by the same model, and the pattern
across them is consistent: it follows a criterion it can evaluate against the
text in front of it, and ignores a disposition about how to feel when unsure.

\paragraph{A shared accelerator breaks reproducibility.} Two test suites run
concurrently produced different verdicts on byte-identical input. The cause is
request batching in the serving daemon. We first misdiagnosed this as model
flakiness and then as a prompt defect, and ``fixed'' it by adding an example,
which made it worse. Determinism required a fixed seed and a machine-wide lock.

\paragraph{An A/B harness must pin everything the code reads.} Our first
before/after comparison reported three regressions. Both arms had been given
different working directories, so one checker saw files on one side only. The
regressions vanished once the directory was pinned.

\section{Threats to validity}

\paragraph{One user, one idiom.} The corpus comes from a single developer. The
patterns carry that developer's phrasing, and the judge required an explicit
reading list for River Plate Spanish because \emph{quedó cableado} (``it is
wired end to end'') was being read as work in progress. Nothing here has been
tested on a second user's transcripts.

\paragraph{The gate contaminates its own corpus.} Gated transcripts were
produced with the gate running, so their closings are partly the gate's own
output distribution. This is why all label-based evaluation is restricted to
ungated transcripts, and why the era comparison in Section~\ref{sec:results} is
observational rather than controlled.

\paragraph{The push label is one-sided.} Absence of a push is not evidence the
turn was complete: a developer often absorbs a premature stop and simply
answers. Every precision in this paper is therefore a floor on the true value
and not an estimate of it. What the label supports is the comparison of rates,
and the comparison holds under any one-sided correction: the gate fires several
times more often than a push occurs.

\paragraph{The grading rule rewards motion.} Counting tool calls after a block
cannot see whether the work was any good. An agent that learned to emit cheap
tool calls after every block would defeat it. We have not observed this, and it
is the most obvious way the loop could be gamed.

\paragraph{Rules are searched on the rows they are scored on.} Every class
demotion, and the three rules added from the missed pushes, were chosen by
reading this corpus and reported on it. Section~
ef{sec:results} splits the
corpus by session and scores the search on held-out sessions to bound the
damage: the direction transfers, the size is about two closings, and the set of
classes chosen is not stable across halves.

\paragraph{The corpus is small where it counts.} 747 closings carry 78 pushes,
so every precision here has a Wilson interval two to four points wide and every
per-class number rests on tens of rows. No configuration we measured separates
from the base rate at 95\%. Conclusions in this paper are stated at the level
the intervals support, which is rates and directions rather than values.

\paragraph{Thresholds are tuned, not derived.} The demotion threshold and the
work threshold were chosen for the volume this log receives.

\section{Discussion}

Three things generalize beyond this system.

\paragraph{Put the decision where the capability is.} The clearest single result
is the background-work case: an exception that a 9B model could not hold in its
prompt became trivial once it was decided by ten lines of transcript inspection.
The temptation with a model in the loop is to route every judgement through it.
The cheaper lanes should take everything they can prove.

\paragraph{Confidence is a router before it is a decision.} Per-pattern
confidence is what makes the local model affordable, because it decides which
messages are worth a model at all. The judge's own confidence, by contrast,
turned out to be publishable and not actionable. A signal can be real, be
logged, be informative to a maintainer, and still not be safe to threshold.

\paragraph{Interventions leave labels.} Any system that interrupts an agent gets
a natural experiment for free: the agent's behaviour after the interruption. We
suspect this generalizes to most guardrails. The evaluation problem for
interventions is often not a missing label but an unread one.

\section{Limitations and future work}

The loop corrects unevenly. It finds patterns that block too much cheaply,
because every block leaves a graded trace and the report reads itself. The
other side needs the push label: a turn the gate allowed and the developer then
pushed back on is a miss, and reading the twelve of them took an afternoon and
produced three rules. That path is manual, and it is only trustworthy on
ungated transcripts, where the label is not partly the gate's own output. On
this corpus the label does not thin out under the gate - 10.3\% of gated
closings carry a push against 10.5\% of ungated ones - so the constraint is
contamination rather than sparsity, and it is the same constraint either way:
the misses can only be read where the gate was not running.

Three extensions follow directly. Demotion is currently proposed and applied by
an agent reading the report; promotion in the other direction is not
implemented. The judge is a general instruction-tuned model, and the graded log
is exactly the dataset needed to fine-tune a specialist. And no part of this has
been run against a second user, which is the first thing that should happen.

\section{Conclusion}

A gate that stops an agent from stopping early is easy to build and hard to
evaluate, because evaluation needs labels and labels need a reader. Both
labels this system needs were already lying in the transcript: the developer's
own reply before the gate exists, and the agent's own behaviour after the gate
intervenes. Neither costs an annotation.

Reading them was uncomfortable and useful in the same measure. The gate fires
on three closings in four and is right about one in eight, which is the rate at
which the failure occurs at all: no configuration we measured separates from it
at 95\%. The first thing the labels
overturned was the design's most confident exemption, which turned out to be
releasing turns worse than average. The second was a prompt written as a
disposition, replaced by one written as a test, which improved every number at
once.

Neither correction came from a new idea. Both came from finally measuring
something that had been sitting in the transcript the whole time.
